% !TEX root = userguide.tex

\chapter{Multiple element groups}
\pgst{tfemch}

\def\chaptertitle{Multiple element groups}
\def\headdate{23rd September 2010}

In this chapter we will briefly describe how to work with multiple element
groups.

\section{Introduction}

The mesh can be splitted into multiple groups, where each group represents a set
of elements. The splitting must already be done at the mesh generation level.
In the basic meshgenerator of \tfem\ this can be achieved using the subroutine
\texttt{mesh\_merge} of the module \texttt{meshgen\_extra}
(see Sec.~\ref{sec:meshgenextra}).

The reasons for introducing multiple groups include
\begin{itemize}
   \item different element shapes, for example triangles and quads in one mesh. 
   \item different materials, material properties or, more generally,
         different structure coefficients in one problem. For example, a Stokes
         problem with two separate regions with a different viscosity.
   \item different physical or mathematical models in separate regions.
         For example with fluid-solid interactions, the fluid and the
         solid will have different element groups.
   \item different physical or mathematical models in partly overlapping
         regions, for example a fluid domain with a ferrofluid influenced by a
         magnetic field. The magnetic field is solved in a larger part of the
         domain than the fluid.
\end{itemize}

\section{Basic usage}

All the routines for assembling the matrix and vector, computing derive vectors,
etc., but also the plotting routines have heading parameters \texttt{groups},
\texttt{elgroups}, \texttt{elgroup2} to perform the operation for
certain groups (e.g.\ \texttt{groups=[2,4]}),
range of groups (e.g.\ \texttt{elgroup1=2}, \texttt{elgroup2=4} for range
2,3,4) or a single group (e.g.\ \texttt{elgroup1=2})
only\footnote{Don't forget to use the
heading parameters \texttt{addmatvec}, \texttt{addvec}, \texttt{add} or some
other variant if you split up an operation over multiple groups. If you do not
use these parameters the output array/scalar will be set to zero at the start
of a routine, destroying the results of a former part of the operation.}. In
that way the assembling process, for example, can be splitted over several
calls of \texttt{build\_system}\footnote{If
element routines are designed to deal with multiple element types, it may not
be necessary to split.}.

\section{Degrees of freedom in connecting nodes of directly connected groups}

If groups are directly connected, i.e.\ there are nodes which belong to more
than one group, the number of degrees of freedom in the system vector
(of type \texttt{sysvector\_t} and
the number of degrees of all the vectors (of type \texttt{vector\_t}) need to
be \emph{the same in common nodes of all connecting groups}. If not, an error
message is produced and the program stops\footnote{If you \emph{really} want
to get beyond this error message you can set
\texttt{allow\_different\_numdegfd\_in\_nodes=.true.}\ before entering 
\texttt{problem\_definition}, but then you are entering messy territory and you
really have to know what you are doing. Lots of pitfalls here. Don't say you
haven't been warned!}. 

The best way to deal with a varying number of degrees between groups is to
disconnect
the element groups (overlapping curves/surfaces are separate and have different nodes).
The connections between the degrees of freedom in the groups can be made using
constraints (see Chapter~\ref{ch:constraints})
on the interface between the groups.

\section{Inactive groups in the problem definition}
\label{sec:inactivegroups}
Inactive groups in the problem definition\footnote{Inactive groups in the mesh
have not been implemented. It could resolve some of the problems
related to inactive groups in problems that use a mesh with connected groups.
However it also makes everything more complicated, since then you have to deal
with both multiple meshes and multiple problems.}
 are element groups that 
\begin{enumerate}
   \item have zero degrees of freedom in the system vector and all the vectors
         (\texttt{input\_probdef\%elnumdegfd=0} and 
          \texttt{input\_probdef\%vec\_elnumdegfd=0}),
   \item are excluded from the build process, derive vector computation,
         plotting etc.
\end{enumerate}
In order to facilitate this without changing a lot of heading parameters and
add the \texttt{groups} parameter, a parameter \texttt{numinactivegroups}
has been added to the
routine \texttt{create\_input\_probdef}. In the following excerpt from a main
program we illustrate the use of this parameter
\begin{listing}{1}
! problem definition

  call create_input_probdef ( mesh, input_probdef, nvec=3, nphysq=2, &
    numinactivegroups=1 )

  input_probdef%vec_elementdof(1)%a =   &
      reshape ( [2,2,2,2,2,2,2,2,2,    &  ! velocity
                 1,0,1,0,1,0,1,0,0,    &  ! pressure
                 1,1,1,1,1,1,1,1,1 ], &  ! scalar, such as vorticity
                 [9,3] )
  input_probdef%vec_elementdof(2)%a = 0

  input_probdef%physq = [1,2]
  input_probdef%inactivegroups = [2]

  ... other statements, like setting essential nodes

  call problem_definition ( input_probdef, mesh, problem )

\end{listing}
Here it is assumed that the mesh has two groups, where group 2 is inactive in
the problem definition.

If you use inactive groups together with groups that have common nodes, an
error will be produced as in the previous section: it is not allowed to have
nodes that have different degrees of freedom in different elements connected to
that node. You can circumvent the error in the same way as suggested in the
previous section, but you are entering messy territory with lots of pitfalls.
A more straightforwards solution, i.e.\ groups with essential degrees only,
is given in the next section.

\section{Element groups having essential degrees of freedom only}

In the routine \texttt{define\_essential} it is possible to set complete
element groups to have essential degrees of freedom only. This maybe useful for
various applications. For example, it can be used as an alternative for
inactive groups, as described in Sec.~\ref{sec:inactivegroups}. In
particular for groups that are connected, it is an attractive approach. Some
notes:
\begin{itemize}
   \item By default all the nodes are set to essential. Use
         \texttt{excludecurves} and/or \texttt{excludesurfaces} to exclude
         some nodes of the group, in particular on the connecting curves or
         surfaces.
   \item If connecting nodes have degrees of freedom that are not
         essential 
         and the group is to be discarded (inactive), the system
         matrix of the
         group should be build with \texttt{zeromatvec=.true.} in
         \texttt{build\_system} to avoid any interactions.
   \item Exclude the inactive groups from computing derivatives arrays
         in \texttt{derive\_vector} when the result is an average of different
         element values of all elements connected, such as gradients.
\end{itemize}
